Full local top-$c$ retention establishes correctness, but can produce
$\gscan c$ intermediate survivors for only $c$ global outputs. The shared-bound
component reduces this excess selection work. The scheduler
then chooses how much scan parallelism to expose. Both operate at a fixed
refinement budget; choosing that budget is the separate problem in
\Cref{ch:refinement}.

\subsection{A complete-key upper bound with explicit witnesses}

Let $K_{(c)}$ be the global $c$th key when $|U|\geq c$. If a group holds
$c$ distinct, fully scored candidates, the largest key among those
witnesses, $\tau_j$, satisfies $K_{(c)}\preceq\tau_j$. There are already
$c$ candidates no later than $\tau_j$ in the global universe. The shared
bound is initialised to a top sentinel and updated by
\begin{equation}
 \tau\leftarrow\min_{\prec}(\tau,\tau_j).
 \label{eq:5-shared-update}
\end{equation}
No bound is published from fewer than $c$ valid candidates. If the entire
universe has fewer than $c$, no group can publish such a bound and all
valid candidates remain eligible.

\begin{resultblock}
\begin{proposition}[Safe rejection by a witnessed bound]
\label{prop:5-shared-bound}
Assume fixed keys, publications backed by $c$ distinct candidates, and
reads that return either the initial sentinel or one complete published
key. Rejecting a candidate only when $K_q(x)\succ\tau$ cannot remove a
global top-$c$ candidate, irrespective of publication timing.
\end{proposition}
\begin{proof}
A publication has a set $W_\tau$ of $c$ distinct witnesses:
\[
 \max_{x\in W_\tau}K_q(x)\preceq\tau
 \quad\Longrightarrow\quad K_{(c)}\preceq\tau.
\]
Sentinel reads and minima of published keys obey the same inequality. Thus
\[
 x\in\operatorname{Top}_c(U)
 \ \Longrightarrow\ K_q(x)\preceq K_{(c)}\preceq\tau_{\rm seen}
 \ \Longrightarrow\ x\text{ is not rejected}.
\]
Stale reads only loosen the bound. Publications
made after earlier safe rejections remain valid: any $c$ retained
distinct candidates still supply $c$ witnesses, whether or not they are
the best $c$ of their group.
\end{proof}
\end{resultblock}

A shared bound may discard local winners while preserving \cref{prop:5-local-coverage}'s global coverage. Admission is
\begin{equation}
 K_q(x)\prec\lambda_j
 \quad\text{and}\quad
 K_q(x)\preceq\tau_{\rm seen},
 \label{eq:5-admission}
\end{equation}
Here $\lambda_j$ is the local compaction threshold and $\tau_{\rm seen}$ the last shared key read. Equality with the shared bound stays eligible.

\subsection{Publication, filtering and implementation limits}

For nonnegative FP32 scores, canonical positive zero and valid nonnegative positions,
\begin{equation}
 \kappa(D,p)=2^{32}\operatorname{uintbits}_{32}(D)+p.
 \label{eq:5-packed-key}
\end{equation}
Non-NaN float bits are increasing on this domain, so
\[
 (D_1,p_1)\prec(D_2,p_2)\iff\kappa(D_1,p_1)<\kappa(D_2,p_2).
\]
Infinity follows finite scores; the all-one sentinel follows valid keys. Invalid positions cannot publish.

\begin{figure}[tbp]
\centering
% Algorithm diagram; the whole-key-read premise is discussed in Section 5.5.
\begin{tikzpicture}[
  font=\small\linespread{1}\selectfont,
  box/.style={draw=ThesisNavy!65,rounded corners=2pt,align=center,
    inner sep=5pt,text width=6.2cm,minimum height=.9cm},
  side/.style={box,text width=4.1cm,fill=ThesisNavy!5},
  arr/.style={-{Latex[length=2mm]},draw=ThesisNavy!80,thick}]
\node[box] (score) at (0,0) {Read primary code; compute full $\primarydist$\\
  Form $K=(\primarydist,p)$};
\node[box,fill=ThesisNavy!5] (admit) at (0,-1.8)
  {Admission: $K\prec\lambda_j$ and $K\preceq\tau_{\rm seen}$};
\node[box] (compact) at (0,-3.8) {Append admitted keys\\
  Compact when necessary\\Keep at most $c$};
\node[box] (trim) at (0,-5.9) {Final compaction and prefix trim\\
  Emit valid keys $K\preceq\tau_{\rm seen}$ to merge};
\node[side] (shared) at (6.3,-.1) {Per-query shared word $Q$\\
  Encodes bound $\tau$\\Initially all ones\\$\tau\succeq K_{(c)}$};
\node[side] (publish) at (6.3,-3.8) {Only with $c$ valid witnesses:\\
  $\operatorname{atomicMin}(Q,\kappa(\tau_j))$};
\draw[arr] (score) -- (admit);
\draw[arr] (admit) -- node[right,font=\footnotesize] {yes} (compact);
\draw[arr] (compact) -- (trim);
\draw[arr] (compact.east) -- (publish.west);
\draw[arr] (publish.north) -- (shared.south);
\draw[arr,dashed] (shared.west) -- ++(-.35,0) |- (admit.east);
\node[side,minimum height=1.1cm] (discard) at (6.3,-6)
  {Discard candidate\\Scoring complete};
\draw[arr] (admit.west) -- node[above,font=\footnotesize] {no} ++(-.5,0) |- ([yshift=-.9cm]trim.south)
  -| (discard.south);
\end{tikzpicture}

\caption[Shared complete-key filtering]{Shared complete-key filtering for one scan group. A candidate's
full primary score is computed before admission. A valid $c$-candidate
witness set can tighten the per-query bound; a stale but complete bound
only weakens filtering. Final prefix trimming reduces the valid merge
input, while the allocated workspace capacity remains unchanged.
The whole-key-read premise is qualified in the text.}
\label{fig:5-shared-bound}
\end{figure}

One thread polls per chunk. After compaction, $c$ valid keys permit a 64-bit \texttt{atomicMin} publication. Final publication/binary search trims the sorted prefix; exact merging handles stale-bound extras. Searches reset bounds in-stream and merge after scanning, without inter-group waiting.

Aligned \path{ld.volatile.global.u64}/\path{atom.global.min.u64} satisfy PTX single-copy atomicity for compatible strong accesses \cite{ptx_isa}. CUDA 13.0 \texttt{ptxas}, \texttt{sm\_120} corroborates this (\cref{sec:a-supplement}); driver-JIT SASS was not captured. C++ \texttt{volatile} alone supplies no such argument \cite{cuda_programming_guide}.

The bound follows full scoring (\cref{fig:5-shared-bound}), saving buffer/merge work while adding polling/atomics. It can slow search (\cref{sec:7-4}).

\subsection{Choosing parallelism without changing the budget}

Grouping launches $b\gscan$ blocks, adding parallelism, staging, selection and merging. Smaller quotas need a certificate/fallback. Concentrated pilot survivors motivated witnessed bounds (\cref{sec:7-4}).

Let $T_{\rm blk}$ be scan threads per block, $H$ the number of SMs times
the maximum resident threads per SM, and
$v=\operatorname{sd}(n_c)/\operatorname{mean}(n_c)$ the candidate-count
coefficient of variation. Its rule is
\begin{align}
 g_{\rm fill}
 &=\max\!\left(\{1\}\cup
       \{h\in\{2,4,8\}:bhT_{\rm blk}\leq0.6H\}\right),
       \label{eq:5-fill-rule}\\
 g_{\rm skew}&=
   \begin{cases}4,&v\geq0.3,\\1,&v<0.3,\end{cases}
 \qquad
 g_{\rm chosen}=\max(g_{\rm fill},g_{\rm skew}).
 \label{eq:5-skew-rule}
\end{align}
The skew rule may override the fill limit. In the evaluation harness,
$v$ is calibrated using 512 database vectors as pseudoqueries and a
fixed routing configuration; it uses no ground-truth labels. This is
distinct from Chapter 6's sampled output-agreement calibration. Once
calibrated, scheduling uses the batch size and stored statistic without
an additional per-search host readback. The constants were fitted on
four pilot configurations and then evaluated on further configurations.
The one-time pseudoquery calibration is excluded from
the reported steady-state search times. Observed startup ranges, separate from refinement-budget calibration (\cref{sec:a-supplement}), are
\[
 T_{\rm gather}\in[0.81,7.04]\,\mathrm{ms},\quad
 T_{\rm first}\in[6.00,27.77]\,\mathrm{ms},\quad
 \operatorname{median}(T_{\rm repeat})\in[2.72,18.13]\,\mathrm{ms}.
\]

Choosing one group dispatches $\gscan=0$ at fixed index/probes/budget. Its distance-only ties need a common-order check for identical survivors. Integrated gains/regret are in \cref{sec:7-6}.

\subsection{Observed execution paths}

On the 170-SM RTX 5090 (48 warps/SM) \cite{blackwell_tuning},
\[
 1000\cdot2\cdot512=1{,}024{,}000
 >0.6\cdot170\cdot48\cdot32=156{,}672.
\]
Thus batch 1,000 selects one group by fill; skew may force four.

Nine of 38 E1 records cross $v=0.3$, all Vogue ($0.338$--$0.622$). Across 132 run--budget cells, Vogue uses four groups in 27; arXiv/OpenAI use $\gscan=0$ in 105. Their full-batch gains use LUTs, packed access, lane refinement and probe selection, without grouped retention/shared bounds.

Vogue has 3--15 tie-only primary differences per 1,000 queries, zero non-tie differences and bit-identical final outputs. The 7--11 range is specific to 128 probes.

Observed cell counts are
\[
 \begin{array}{c|rrr}
 &G=0&G=2&G=4\\\hline
 \text{held-out}&10&0&2\\
 S=128&2&8&2
 \end{array}
\]
Eight of twelve switch paths. Grouping serves smaller/skewed loads, without a fixed batch-256 boundary.

The heuristic omits register/shared residency and compaction/merge costs; pseudoqueries may age. Worst recorded regret is 1.217 at batch 256. These tests hold the budget fixed; \cref{sec:5-6} refines survivors.